\documentclass[11pt]{article}
\usepackage{fontspec}
\setmainfont{Times New Roman}
\setsansfont{Arial}
\setmonofont{Consolas}
\usepackage{amsmath,amssymb,mathtools}
\usepackage{geometry}
\usepackage{microtype}
\newcommand{\mo}{\mathfrak{o}}
\newcommand{\mO}{\mathfrak{O}}
\newcommand{\mT}{\mathfrak{T}}
\geometry{a4paper,margin=2.35cm}
\setlength{\parindent}{1.25em}
\setlength{\parskip}{0pt}
\makeatletter
\renewcommand\footnotesize{\@setfontsize\footnotesize{7.7}{8.7}\selectfont}
\makeatother
\emergencystretch=100em
\allowdisplaybreaks
\sloppy
\hbadness=10000
\vbadness=10000
\hfuzz=2pt
\vfuzz=10pt
\raggedbottom

\begin{document}

% Unit: Paper31 Section07 no. 1 opening v001. Direct Interslavic Latin translation from the German final-audited source slice, source lines 385--389 / segments P31-S0082--P31-S0084, expanded against printed scan/OCR witness pages 98--99.
\setcounter{footnote}{28}

\subsection*{§7. Teorem o diskriminantu za porjadky algebraičnogo čislovogo polja ili funkcionalnogo polja}

\noindent\textbf{1.}\quad Razsmatrivane čislove i funkcionalne polja -- pri čem pri algebraičnyh funkcijah od več než jednej neoznačene i pri cělyh algebraičnyh funkcijah ide se o razširenje funkcionalnoj oblasti -- vse podlagajut se slědujučemu tipu polja s fiksovanym pojmom cělyh elementov. Nehaj $\mo$ bude kolco glavnyh idealov bez děliteljev nuly i s jediniceju (kolco cělyh racionalnyh čisel, oblast polinomov od jednej neoznačene s koeficientami iz polja, funkcionalna oblast), a $\Omega$ jego polje kvocientov. Nehaj $K$ bude konečno razširenje prvogo roda polja $\Omega$, a $\mO$ kolco vsih elementov iz $K$, cělyh vzhodno k $\mo$. Vsako podkolco od $\mO$, ktoro obejmaje $\mo$, nazyva se porjadok; samo $\mO$ označaje se kako glavny porjadok.

Vsaky $\mo$-modul v $\mO$, osoblivo vse idealy porjadka i sam porjadok, imaje linearno nezavisnu modulnu bazu vzhodno k $\mo$. Ako $K$ imaje stepen $n$ nad $\Omega$, togda dosta je ograničiti se na porjadky ranga $n$ vzhodno k $\mo$, bo vsaky porjadok nižšego ranga črez tvorjenje kvocientov poraždaje medžupolje medžu $K$ i $\Omega$ imenno toj stepeni, i za jego vse predposylky ostavajut. Zato vsaky porjadok $\mT$ jest kolco, za ktoro predposylky §4 i §5 sut ispolnjene, i imenno kolco bez děliteljev nuly; tym diskriminant $D_\mT$ porjadka jest jednoznačno opreděljen do faktora, ktory jest kvadrat jedinice iz $\mo$. Ta diskriminant jest takože, do faktora-jedinice iz $\Omega$, identičny s diskriminantom polja $K$, smatrjanym kako $\Omega$-modul, i zato različny od nuly, bo $K$ byl predpoloženy kako razširenje prvogo roda (§6, 3). Diskriminant takože dopušča predstavjenje, dano v §6, 4, črez kvadrat determinanta konjugovanyh baznyh elementov.

\end{document}
